% Paper 6 (GRPO-Registry) worked examples
\section{Worked Examples}
\label{sec:p6-worked}
Both examples below are verbatim CLI runs on the shipped entries.

\subsection{Example 1: \texttt{stackdiff} on the ``DAPO'' Pair}
\label{sec:p6-worked-dapo}
The registry contains two stacks claiming the label \texttt{dapo}: the open
Colab trainer arm and the managed-runtime arm. Diffing them:

\begin{lstlisting}[basicstyle=\ttfamily\scriptsize, frame=single,
  breaklines=true, caption={\texttt{stackdiff colab-open_dapo_e3
  tinker_dapo_qwen3.5-4b_gsm8k} (abridged to 10 of 37 output lines). Regenerate with \texttt{python3 registry/query.py stackdiff}.},
  label={lst:p6-stackdiff}]
stackdiff colab-open_dapo_e3 vs tinker_dapo_qwen3.5-4b_gsm8k
  shared claimed label: 'dapo'
  [R5] checkpoint.id: 'Qwen/Qwen2.5-0.5B-Instruct' -> 'Qwen/Qwen3.5-4B'
  [R5] data.dataset: 'synthetic a+b, a,b in [11,60]' -> 'openai/gsm8k:main'
  [R5] eval.metric: 'greedy held-out accuracy delta (post - pre)' -> 'mean sampled train reward over last 10 steps'
  [R4] loss.objective: 'ppo-clip on sequence log-prob ratio, summed over sequences' -> 'reinforce: -A_i * sum_t logp, mean over sequences'
  [R4] loss.clip: 'ratio[0.8,1.28]' -> 'advantage[-2,3]'
  [R4] loss.dynamic_sampling: 'resample zero-variance groups up to 2 retries, then drop' -> False
verdict: R5 outcome-critical differences (checkpoint, reward, prompt/chat template, dataset, eval protocol, execution mode)
coverage: 0.87 of compared fields known on both sides; 7 unknown (max hidden severity R5)
\end{lstlisting}

The verdict is R5: the two ``DAPO'' entries differ in checkpoint, task,
reward, prompt template, and evaluation protocol, before any algorithmic
field is considered. Among the R4 differences, dynamic sampling---the component
that mechanically drives batch ZVF to zero---is present on the open side and
absent on the managed side, whose ``dapo'' arm is a REINFORCE-style loss with
an advantage clip. This is a
\emph{pre-run} verdict computed from manifests alone. The two arms also differ
in model, task, $G$, steps, and seeds (all R5/R3 differences), so their ZVF
values ($0.00$ vs.\ $0.58$) are not a like-for-like comparison, and the
managed ``dapo'' arm is an advantage-clipping variant rather than DAPO; the
defensible telemetry is within-stack (GRPO$\to$``DAPO'': $0.25\to0.00$ open,
$0.567\to0.578$ managed). The general lesson matches the under-specified
same-label exhibit, where checkpoint, reward, prompt, batch, data, and loss
fields co-varied as final reward spanned $85.6\% \to 5.0\%$: rankings
inherit the stack, and \texttt{stackdiff} makes the inheritance legible
before compute is spent.

\subsection{Controlled Validation: Matched TRL-vs-Tinker Evaluation}
\label{sec:p6s-claimc-slot}
To verify that \texttt{stackdiff} correctly isolates confounds, we executed a
controlled head-to-head experiment between TRL and Tinker
(\texttt{experiments/results/publication\_fixes/matched\_trl\_tinker/}).
Holding all R5 and R4/R3 fields fixed---identical base checkpoint (\texttt{Qwen/Qwen3-8B}),
prompt template, regex answer extractor, decoding parameters, batch shape
($16$ prompts/step, $G{=}8$), and GSM8K data split across three seeds ($42, 123, 456$)
for 30 steps---reduces the \texttt{stackdiff} verdict to R2 (runtime differences
only) at high coverage.

Under this matched configuration, the large nominal gap disappears into noise:
(i) the last-$10$ training reward difference (Tinker minus TRL) has mean $+0.050$
($95\%$ $t$-CI $[-0.488,\,+0.588]$, paired bootstrap $[-0.156,\,+0.294]$);
(ii) held-out greedy evaluation on $200$ test items shows a mean accuracy difference
of $+0.012$ ($95\%$ two-way bootstrap CI $[-0.050,\,+0.078]$, $t$-CI $[-0.109,\,+0.132]$).
Both confidence intervals span zero. This empirical validation confirms that the
$17\times$ nominal gap was entirely an artifact of unrecorded R5/R3 confounds,
and that the R0--R5 ladder successfully identifies true comparability boundaries.


\subsection{Example 2: ``Which Stacks Report KL Handling?''}
\label{sec:p6-worked-kl}
Item 2 (reference policy + KL) is the canonical scattered-in-appendices
field: defaults differ silently across frameworks ($\beta{=}0.04$ in TRL and
veRL versus $0.02$ in OpenRLHF, per their own dumps), and DAPO removes the
term outright. The registry answers in one call:

\begin{lstlisting}[basicstyle=\ttfamily\scriptsize, frame=single,
  caption={\texttt{query --item reference_kl} on the seed entries
  (grouped).}, label={lst:p6-klquery}]
colab-open_{grpo,drgrpo,dapo,grpo-adaptiveg}   reference_kl=FULL (1.00)
trl_grpo_qwen3-8b_gsm8k                        reference_kl=FULL (1.00)
verl_grpo_qwen3-8b_gsm8k                       reference_kl=FULL (1.00)
openrlhf_grpo_qwen3-8b_gsm8k                   reference_kl=FULL (1.00)
tinker_{grpo(8B),grpo,dapo,drgrpo,gspo}        reference_kl=PARTIAL (0.33)
\end{lstlisting}

Seven of twelve stacks fully report KL handling. The instructive cases are the
two ends: the Colab arms score \texttt{FULL} by \emph{explicitly declaring the
absence} of a reference policy ($\beta{=}0.0$, estimator \texttt{none})---a
reported negative is a first-class answer---while all five managed-runtime
entries are \texttt{PARTIAL} because \texttt{kl.beta} and the estimator are
\texttt{managed_by_tinker} nulls. A meta-analyst asking ``does KL
regularization explain the variant ranking?'' learns, in one query, exactly
which stacks can enter the regression and which must be excluded as
unauditable on that axis.
